\documentclass[12pt,a4paper]{article}
% Globalni nastaveni dokumentu.
\usepackage[a4paper,left=3.2cm,right=3.2cm,top=2.7cm,bottom=2.9cm,headheight=15pt]{geometry}
\usepackage{fontspec}
\usepackage{polyglossia}
\setdefaultlanguage{czech}
\setmainfont{TeX Gyre Pagella}
\setsansfont{TeX Gyre Heros}
\usepackage{microtype}
\usepackage{enumitem}
\usepackage{hyperref}
\usepackage{xcolor}
\usepackage{amsmath,amssymb}
\DeclareMathOperator*{\plim}{plim}
\usepackage{titlesec}
\usepackage{tocloft}
\setlength{\cftsecnumwidth}{2.6em}
\usepackage{fancyhdr}
\hypersetup{colorlinks=true,linkcolor=black,urlcolor=blue}
\linespread{1.08}
\setlength{\parindent}{0pt}
\setlength{\parskip}{0.45\baselineskip plus 1pt}
\setlength{\emergencystretch}{2em}

\makeatletter
\renewcommand*\l@section{\@dottedtocline{1}{0em}{2.4em}}
\makeatother

\setlist[itemize]{topsep=4pt,itemsep=2pt,parsep=0pt,leftmargin=1.4em}
\setlist[enumerate]{topsep=4pt,itemsep=2pt,parsep=0pt,leftmargin=1.7em}
\titleformat{\section}{\Large\bfseries\sffamily}{\thesection}{0.6em}{}
\titleformat{\subsection}{\large\bfseries\sffamily}{\thesubsection}{0.6em}{}
\titleformat{\subsubsection}{\normalsize\bfseries\sffamily}{\thesubsubsection}{0.6em}{}
\titlespacing*{\section}{0pt}{1.2\baselineskip}{0.55\baselineskip}
\titlespacing*{\subsection}{0pt}{0.9\baselineskip}{0.35\baselineskip}
\titlespacing*{\subsubsection}{0pt}{0.7\baselineskip}{0.25\baselineskip}
\pagestyle{fancy}
\fancyhf{}
\lhead{\small Příprava ke státnicím - Data Analytics}
\rhead{\small\thepage}
\cfoot{}
\newcommand{\term}[1]{\textbf{#1}}
\newcommand{\uv}[1]{„#1“}
\newcommand{\todohead}{\subsection*{Co doplnit mimo dodané materiály}}
\newcommand{\pozor}{\subsection*{Pozor u zkoušky}}


\newenvironment{terminy}
  {\begin{itemize}[leftmargin=1.4em,itemsep=4pt]}
  {\end{itemize}}
\newcommand{\defterm}[2]{\item \term{#1}: #2}

\begin{document}

\begin{titlepage}
  \centering
  \vspace*{2.4cm}
  {\Large\bfseries Bakalářské studium\par}
  \vspace{1.2cm}
  {\Huge\bfseries Shrnutí otázek a klíčové termíny\par}
  \vspace{0.6cm}
  {\LARGE\bfseries Data Analytics\par}
  \vspace{0.8cm}
  \rule{0.72\textwidth}{0.4pt}\par
  \vspace{0.8cm}
  {\large Rychlý opakovací přehled ke státní závěrečné zkoušce\par}
  \vfill
\end{titlepage}

\tableofcontents
\newpage

\section{Transakční data}

\subsection{Shrnutí}
Transakční data zachycují konkrétní provozní události organizace ve vysoké granularitě, například
objednávky, faktury, platby, skladové pohyby nebo servisní požadavky. Vznikají hlavně v OLTP
systémech, které jsou optimalizované na rychlý a konzistentní zápis jednotlivých změn. Pro analytiku
se tato data obvykle přesouvají datovými pumpami do DSA, ODS, DWH, datových tržišť nebo OLAP vrstev.
Důležité je rozlišit ETL a ELT, dávkové a průběžné zpracování, CDC, streamy a fronty zpráv. Při
návrhu pipeline se kromě samotného přesunu řeší kvalita dat, idempotence, auditovatelnost,
bezpečnost, monitoring a latence.

\subsection{Nejdůležitější termíny}
\begin{terminy}
  \defterm{Transakční data}{Záznamy o konkrétních provozních událostech, typicky na úrovni dokladu,
  položky, platby, skladového pohybu nebo změny stavu.}
  \defterm{Byznys transakce}{Událost v podniku, která mění stav byznysu, například vystavení faktury
  nebo přijetí objednávky.}
  \defterm{Databázová transakce}{Technická jednotka práce v databázi, která má proběhnout celá, nebo
  vůbec.}
  \defterm{ACID}{Soubor vlastností databázové transakce: atomicita, konzistence, izolovanost a
  trvanlivost.}
  \defterm{OLTP}{Provozní zpracování transakcí optimalizované na velké množství krátkých zápisů,
  změn a čtení jednotlivých záznamů.}
  \defterm{ETL}{Proces extract-transform-load, kdy se data nejprve načtou, potom transformují a až
  poté uloží do cílového analytického úložiště.}
  \defterm{ELT}{Proces extract-load-transform, kdy se data nejprve uloží do cílové platformy a
  transformace probíhají až v ní.}
  \defterm{DSA}{Data staging area; dočasná technická vrstva pro extrahovaná data před jejich
  vyčištěním a konsolidací.}
  \defterm{ODS}{Operational data store; konsolidované úložiště aktuálních provozních dat pro
  operativní reporting a integraci.}
  \defterm{DWH}{Datový sklad; integrované, historické, subjektově orientované a relativně neměnné
  úložiště pro podporu rozhodování.}
  \defterm{CDC}{Change data capture; zachytávání a přenos pouze změněných, vložených nebo smazaných
  záznamů.}
  \defterm{Idempotence pipeline}{Vlastnost, že opakovaný běh zpracování nevytvoří duplicity ani
  nezmění správný výsledek.}
\end{terminy}

\section{Metodika výběrových šetření}

\subsection{Shrnutí}
Výběrové šetření zkoumá vzorek a výsledky zobecňuje na cílovou populaci. Kvalita výsledku závisí na
jasném cíli, správném vymezení populace, opory výběru, výběrového designu, dotazníku, pilotáže,
sběru, čištění, vážení a interpretace. Pravděpodobnostní výběry umožňují odhad výběrové chyby,
nepravděpodobnostní výběry jsou spíše praktickým nebo explorativním kompromisem. Dotazník má být
jednoznačný, nezaujatý a pilotovaný. Vedle výběrové chyby je nutné hlídat chybu pokrytí, neodpovědi,
měření a zpracování. Vážení upravuje vzorek podle designu, neodpovědí a známé struktury populace.

\subsection{Nejdůležitější termíny}
\begin{terminy}
  \defterm{Cílová populace}{Množina jednotek, o které chce šetření vypovídat.}
  \defterm{Výběrová jednotka}{Jednotka, která může být vybrána do vzorku, například osoba, domácnost,
  firma nebo prodejna.}
  \defterm{Opora výběru}{Seznam nebo mechanismus, z něhož se vybírají jednotky do vzorku.}
  \defterm{Vzorek}{Skutečně vybraná část populace, ze které se sbírají data.}
  \defterm{Respondent}{Jednotka, od které se podařilo získat odpověď.}
  \defterm{Pravděpodobnostní výběr}{Výběr, u kterého má každá jednotka známou nenulovou pravděpodobnost
  zařazení do vzorku.}
  \defterm{Stratifikovaný výběr}{Rozdělení populace do vrstev a výběr jednotek z každé vrstvy; často
  zvyšuje přesnost a zajišťuje zastoupení důležitých skupin.}
  \defterm{Shlukový výběr}{Výběr skupin jednotek, například škol nebo obcí; bývá levnější, ale může
  zvýšit výběrovou chybu.}
  \defterm{Pilotáž}{Ověření dotazníku a sběru dat před ostrým šetřením.}
  \defterm{Neodpověď}{Situace, kdy vybraná jednotka neposkytne odpověď nebo chybí odpověď na konkrétní
  otázku.}
  \defterm{Designová váha}{Váha obvykle rovná převrácené pravděpodobnosti výběru jednotky.}
  \defterm{Raking}{Iterativní úprava vah tak, aby vzorek odpovídal známým marginálním rozdělením
  populace.}
\end{terminy}

\section{Externí datové a informační zdroje}

\subsection{Shrnutí}
Externí zdroje doplňují interní data o tržní, regulatorní, makroekonomický, regionální nebo
konkurenční kontext. Patří sem oficiální statistiky, registry, otevřená data, komerční databáze,
odborné publikace, webové zdroje, API nebo OSINT. Před použitím je nutné posoudit důvěryhodnost,
metodiku, aktuálnost, granularitu, právní podmínky a technickou dostupnost zdroje. Competitive
intelligence je legální a etická práce s informacemi o vnějším prostředí organizace. Externí data
v pipeline vyžadují metadata, validace, verzování, monitoring, mapování na interní číselníky a
kontrolu licencí.

\subsection{Nejdůležitější termíny}
\begin{terminy}
  \defterm{Externí datový zdroj}{Zdroj dat mimo organizaci, který doplňuje interní pohled o širší
  kontext.}
  \defterm{Primární zdroj}{Původní zdroj dat nebo informací, například registr, měření nebo originální
  dokument.}
  \defterm{Sekundární zdroj}{Zdroj, který původní data shrnuje, interpretuje, katalogizuje nebo dále
  zpracovává.}
  \defterm{Otevřená data}{Data zveřejněná tak, aby byla strojově čitelná, licenčně použitelná a
  dohledatelná.}
  \defterm{OSINT}{Open-source intelligence; získávání a vyhodnocování informací z veřejně dostupných
  zdrojů.}
  \defterm{Competitive intelligence}{Systematický sběr, hodnocení a analýza informací o konkurenci,
  trhu a vnějším prostředí pro podporu rozhodování.}
  \defterm{API}{Programové rozhraní umožňující automatizovaný přístup k datům nebo službám.}
  \defterm{Web scraping}{Automatizované získávání dat z webových stránek, obvykle přes HTML, DOM,
  CSS selektory nebo XPath.}
  \defterm{Metadata}{Data o datech, například původ, metodika, čas platnosti, licence, jednotky,
  granularita a kvalita.}
  \defterm{Granularita}{Úroveň detailu dat, například jednotlivá transakce, den, region nebo firma.}
  \defterm{Licence dat}{Právní podmínky, které určují, zda a jak lze data používat, sdílet nebo
  komerčně využít.}
\end{terminy}

\section{Etika a regulace při sběru a uchovávání dat}

\subsection{Shrnutí}
Etika řeší férovost, přiměřenost, dopady a odpovědnost při práci s daty, zatímco regulace stanoví
závazná pravidla. V evropském prostředí je hlavním rámcem GDPR, které chrání osobní údaje fyzických
osob a vyžaduje právní titul, transparentnost, účelové omezení, minimalizaci, přesnost, omezení
uložení, integritu, důvěrnost a odpovědnost. Zásadní je rozlišit správce, zpracovatele, subjekt
údajů, osobní údaje, zvláštní kategorie údajů, anonymizaci a pseudonymizaci. U datových a
algoritmických systémů je nutné řešit bias, vysvětlitelnost, lidský dohled, auditovatelnost, bezpečnost
a možnost nápravy.

\subsection{Nejdůležitější termíny}
\begin{terminy}
  \defterm{Informační etika}{Oblast etiky zabývající se vznikem, šířením, ukládáním, používáním a
  dopady informací.}
  \defterm{PAPA}{Masonův rámec informační etiky: privacy, accuracy, property a accessibility.}
  \defterm{GDPR}{Obecné nařízení EU o ochraně osobních údajů vztahující se na zpracování osobních
  údajů fyzických osob.}
  \defterm{Osobní údaj}{Jakákoli informace vztahující se k identifikované nebo identifikovatelné
  fyzické osobě.}
  \defterm{Zpracování}{Široký pojem zahrnující sběr, uložení, úpravu, vyhledání, použití, předání
  i výmaz osobních údajů.}
  \defterm{Správce}{Subjekt, který určuje účely a prostředky zpracování osobních údajů.}
  \defterm{Zpracovatel}{Subjekt, který zpracovává osobní údaje pro správce a podle jeho pokynů.}
  \defterm{Subjekt údajů}{Člověk, kterého se osobní údaje týkají.}
  \defterm{Právní titul}{Důvod, který činí zpracování osobních údajů zákonným, například souhlas,
  smlouva, právní povinnost nebo oprávněný zájem.}
  \defterm{Pseudonymizace}{Nahrazení identifikátorů kódem nebo aliasem; identifikace je stále možná
  s dodatečnými informacemi.}
  \defterm{Anonymizace}{Úprava dat tak, že jednotlivce nelze rozumně pravděpodobnými prostředky
  identifikovat.}
  \defterm{DPIA}{Posouzení vlivu na ochranu osobních údajů u rizikového zpracování.}
  \defterm{Fairness modelu}{Požadavek, aby model nepřiměřeně neznevýhodňoval určité skupiny.}
\end{terminy}

\section{Databáze a SQL}

\subsection{Shrnutí}
Databáze je organizovaná množina perzistentních dat spravovaná systémem řízení báze dat. Relační
model ukládá data do tabulek a vztahy vyjadřuje pomocí klíčů. Návrh databáze probíhá na konceptuální,
logické a fyzické úrovni a používá entity, atributy, vztahy, kardinality a integritní omezení.
Normalizace omezuje redundanci a aktualizační anomálie, zatímco denormalizace se používá kontrolovaně
kvůli výkonu nebo analytice. SQL zahrnuje definici struktur, manipulaci s daty, dotazování, řízení
oprávnění a transakce. V praxi je důležité rozumět joinům, agregacím, poddotazům, pohledům, indexům,
exekučním plánům, izolaci transakcí a rozdílu mezi OLTP a analytickým DWH/OLAP návrhem.

\subsection{Nejdůležitější termíny}
\begin{terminy}
  \defterm{Databáze}{Integrovaná a organizovaná množina perzistentních dat.}
  \defterm{SŘBD}{Systém řízení báze dat; software pro správu, dotazování, bezpečnost, transakce,
  zálohování a obnovu databáze.}
  \defterm{Relace}{Tabulka v relačním modelu, tvořená řádky a sloupci.}
  \defterm{Primární klíč}{Vybraný kandidátní klíč, který jednoznačně identifikuje řádek tabulky.}
  \defterm{Cizí klíč}{Atribut nebo skupina atributů odkazující na klíč v jiné nebo stejné tabulce.}
  \defterm{Referenční integrita}{Pravidlo, že cizí klíč musí odkazovat na existující záznam, nebo
  mít povolenou prázdnou hodnotu.}
  \defterm{ER model}{Konceptuální model popisující entity, atributy, vztahy, kardinality a volitelnost.}
  \defterm{Normalizace}{Postup návrhu schématu, který omezuje redundanci a anomálie vkládání,
  aktualizace a mazání.}
  \defterm{3NF}{Třetí normální forma; neklíčové atributy nemají tranzitivně záviset na klíči přes
  jiné neklíčové atributy.}
  \defterm{SQL}{Standardizovaný jazyk pro práci s relačními databázemi.}
  \defterm{JOIN}{Operace spojující řádky z více tabulek podle zadané podmínky.}
  \defterm{Index}{Pomocná datová struktura zrychlující vyhledávání, spojování a řazení, obvykle za
  cenu pomalejších zápisů.}
  \defterm{Exekuční plán}{Popis postupu, kterým databázový systém provede dotaz.}
  \defterm{Deadlock}{Situace, kdy transakce čekají na zámky držené navzájem a nemohou pokračovat.}
\end{terminy}

\section{Vizualizace dat}

\subsection{Shrnutí}
Vizualizace převádí data do vizuální podoby pro analýzu, komunikaci a rozhodování. Explorační
vizualizace pomáhá hledat vzory, chyby a odlehlé hodnoty, zatímco explanatorní vizualizace vysvětluje
konkrétní sdělení publiku. Volba grafu vychází z analytické otázky: trend, porovnání, rozdělení,
vztah, kompozice, geografie, tok nebo nejistota. Dobrá vizualizace má jasný název, popsané osy,
jednotky, zdroj, vhodné měřítko, přiměřené barvy a minimum vizuálního šumu. Dashboard musí podporovat
rozhodování, mít jasné metriky a respektovat roli uživatele. Etická vizualizace nezkresluje data,
neskrývá nejistotu a bere v úvahu přístupnost.

\subsection{Nejdůležitější termíny}
\begin{terminy}
  \defterm{Vizualizace dat}{Převod dat do vizuální podoby pro pochopení vztahů, trendů, rozdílů a
  odchylek.}
  \defterm{Explorační vizualizace}{Pracovní vizualizace používaná analytikem při hledání vzorů,
  chyb a hypotéz.}
  \defterm{Explanatorní vizualizace}{Vysvětlující graf nebo dashboard určený publiku k pochopení
  hlavního sdělení.}
  \defterm{Vizuální kanál}{Vizuální vlastnost použitá ke kódování dat, například poloha, délka,
  barva, tvar nebo velikost.}
  \defterm{Preatentivní vlastnost}{Vizuální znak, který člověk zachytí velmi rychle, například
  výrazná barva nebo izolovaný bod.}
  \defterm{Histogram}{Graf rozdělení metrické proměnné pomocí intervalů.}
  \defterm{Paretův graf}{Kombinace sloupcového grafu seřazených kategorií a kumulativní křivky,
  používaná pro hledání hlavních příčin.}
  \defterm{Boxplot}{Graf shrnující rozdělení pomocí mediánu, kvartilů, IQR, vousů a odlehlých hodnot.}
  \defterm{Sekvenční škála}{Barevná škála pro seřazené hodnoty od nízkých po vysoké.}
  \defterm{Divergentní škála}{Barevná škála pro odchylky kolem významného středu, například kolem
  nuly nebo plánu.}
  \defterm{Dashboard}{Uspořádaná sada vizualizací, ukazatelů a ovládacích prvků pro sledování
  určité oblasti.}
  \defterm{Drill-down}{Přechod z agregovaného pohledu na detailnější úroveň dat.}
\end{terminy}

\section{Deskriptivní statistická analýza}

\subsection{Shrnutí}
Deskriptivní statistika popisuje již pozorovaná data a pomáhá porozumět jejich poloze, variabilitě,
tvaru rozdělení, vztahům a ekonomickým souvislostem. Volba statistiky závisí na typu proměnné:
nominální, ordinální, intervalové nebo poměrové. Základ tvoří četnosti, průměry, medián, kvantily,
rozptyl, směrodatná odchylka, IQR, variační koeficient, histogramy a boxploty. U ekonomických dat je
důležitá agregace podle správné granularity, komparace v čase, indexy, tempa růstu a srovnání s
benchmarkem. Při interpretaci se hlídají chybějící hodnoty, outliery, sezonnost, inflace, změny
metodiky a rozdíl mezi popisem, inferencí, predikcí a kauzalitou.

\subsection{Nejdůležitější termíny}
\begin{terminy}
  \defterm{Deskriptivní analýza}{Popis dat tak, aby bylo zřejmé, co se v nich stalo, bez automatického
  dokazování kauzality.}
  \defterm{Hromadný jev}{Jev vyskytující se u mnoha jednotek nebo opakovaně v čase, takže má smysl
  jej statisticky popisovat.}
  \defterm{Četnost}{Počet nebo podíl výskytů hodnoty či kategorie v datech.}
  \defterm{Průměr}{Míra polohy citlivá na extrémní hodnoty; aritmetický průměr je součet hodnot
  dělený jejich počtem.}
  \defterm{Medián}{Prostřední hodnota po seřazení dat; robustnější vůči extrémům než průměr.}
  \defterm{Modus}{Nejčastější hodnota nebo kategorie.}
  \defterm{Kvantily}{Hodnoty dělící seřazená data podle pořadí; medián je 50\% kvantil.}
  \defterm{Rozptyl}{Průměrná kvadratická odchylka hodnot od průměru.}
  \defterm{Směrodatná odchylka}{Druhá odmocnina rozptylu, vyjádřená ve stejné jednotce jako původní
  proměnná.}
  \defterm{IQR}{Interkvartilové rozpětí \(Q_3-Q_1\), robustní míra variability.}
  \defterm{Variační koeficient}{Relativní míra variability, obvykle poměr směrodatné odchylky a
  průměru.}
  \defterm{Agregace}{Shrnutí detailních dat na vyšší úroveň pomocí funkcí jako součet, průměr,
  počet nebo medián.}
  \defterm{Bazický index}{Index porovnávající všechna období se stejným základním obdobím.}
  \defterm{Řetězový index}{Index porovnávající každé období s obdobím bezprostředně předchozím.}
  \defterm{\(p\)-value}{Pravděpodobnost výsledku alespoň tak extrémního, jako je pozorovaný, za
  předpokladu platnosti nulové hypotézy.}
\end{terminy}

\section{Exploratorní analýza}

\subsection{Shrnutí}
Exploratorní analýza slouží k pochopení dat, odhalení problémů, nalezení vzorů a přípravě na další
modelování. Redukce rozměru zjednodušuje mnohorozměrná data při snaze zachovat podstatnou informaci.
PCA nahrazuje původní kvantitativní proměnné menším počtem nekorelovaných komponent seřazených podle
vysvětlené variability. Faktorová analýza předpokládá latentní faktory vysvětlující korelační strukturu
proměnných. Korespondenční analýza vizualizuje vztahy v kontingenčních tabulkách. Diskriminační
analýza je supervised metoda hledající lineární kombinace proměnných, které oddělují známé skupiny.
Všechny tyto metody jsou citlivé na škálování, chybějící hodnoty, outliery a interpretaci výstupů.

\subsection{Nejdůležitější termíny}
\begin{terminy}
  \defterm{EDA}{Exploratory data analysis; postupy pro počáteční pochopení struktury, kvality a
  vzorů v datech.}
  \defterm{Redukce dimenze}{Nahrazení původních proměnných menším počtem nových proměnných nebo
  reprezentací.}
  \defterm{PCA}{Analýza hlavních komponent; lineární transformace kvantitativních proměnných na
  nekorelované komponenty s maximální vysvětlenou variabilitou.}
  \defterm{Hlavní komponenta}{Lineární kombinace původních proměnných; první komponenta vysvětluje
  největší možnou variabilitu.}
  \defterm{Vlastní číslo}{U PCA vyjadřuje množství variability vysvětlené danou komponentou.}
  \defterm{Faktorová analýza}{Metoda modelující korelační strukturu proměnných pomocí menšího počtu
  latentních faktorů.}
  \defterm{Faktorová zátěž}{Koeficient vztahu mezi původní proměnnou a latentním faktorem.}
  \defterm{Komunalita}{Podíl variability proměnné vysvětlený společnými faktory.}
  \defterm{Rotace faktorů}{Transformace faktorového řešení pro lepší interpretovatelnost, například
  varimax.}
  \defterm{Korespondenční analýza}{Metoda pro analýzu kontingenčních tabulek založená na odchylkách
  od nezávislosti.}
  \defterm{Inerce}{V korespondenční analýze míra variability nebo odchylky od nezávislosti.}
  \defterm{Diskriminační analýza}{Supervised metoda hledající kombinace proměnných, které oddělují
  předem známé skupiny.}
  \defterm{Mahalanobisova vzdálenost}{Vzdálenost zohledňující kovarianční strukturu proměnných.}
\end{terminy}

\section{Lineární regrese}

\subsection{Shrnutí}
Lineární regrese modeluje podmíněnou střední hodnotu závislé proměnné jako lineární kombinaci
vysvětlujících proměnných. OLS odhaduje parametry minimalizací součtu čtverců reziduí. Regrese může
sloužit k deskripci, predikci i kauzální analýze, ale kauzální interpretace vyžaduje především
exogenitu a promyšlený výběr proměnných. Při splnění základních předpokladů je OLS nestranný a
konzistentní, při homoskedasticitě také BLUE. V praxi je nutné řešit interpretaci koeficientů,
diagnostiku reziduí, heteroskedasticitu, multikolinearitu, vynechané proměnné, transformace,
regularizaci, volbu modelu a kvalitu predikce mimo trénovací data.

\subsection{Nejdůležitější termíny}
\begin{terminy}
  \defterm{Lineární regresní model}{Model \(y=X\beta+u\), kde je závislá proměnná vysvětlena lineární
  kombinací prediktorů a chybovou složkou.}
  \defterm{OLS}{Metoda nejmenších čtverců; odhad volící parametry tak, aby minimalizoval součet
  čtverců reziduí.}
  \defterm{Reziduum}{Rozdíl mezi pozorovanou a modelem predikovanou hodnotou.}
  \defterm{Exogenita}{Předpoklad, že chyba má nulovou podmíněnou střední hodnotu vzhledem k
  vysvětlujícím proměnným.}
  \defterm{Homoskedasticita}{Konstantní rozptyl chyb podmíněně na vysvětlujících proměnných.}
  \defterm{Heteroskedasticita}{Nekonstantní rozptyl chyb; ovlivňuje standardní chyby a inferenci.}
  \defterm{Multikolinearita}{Silná lineární závislost mezi vysvětlujícími proměnnými, která zvyšuje
  nejistotu odhadů.}
  \defterm{Gaussova-Markovova věta}{Při splnění podmínek je OLS nejlepší lineární nestranný odhad,
  tedy BLUE.}
  \defterm{\(R^2\)}{Podíl variability závislé proměnné vysvětlený modelem.}
  \defterm{t-test}{Test hypotézy o jednom regresním koeficientu.}
  \defterm{F-test}{Test více omezení nebo společné významnosti skupiny proměnných.}
  \defterm{Omitted variable bias}{Zkreslení odhadu vzniklé vynecháním relevantní proměnné korelované
  se zahrnutým regresorem.}
  \defterm{Ridge}{Lineární regrese s \(L_2\) penalizací velikosti koeficientů.}
  \defterm{LASSO}{Lineární regrese s \(L_1\) penalizací, která může některé koeficienty stáhnout
  přesně na nulu.}
\end{terminy}

\section{Zobecněné lineární regresní modely}

\subsection{Shrnutí}
Zobecněné lineární modely a modely s omezenou závisle proměnnou se používají tehdy, když klasická
lineární regrese neodpovídá povaze \(y\). U binární proměnné se modeluje pravděpodobnost jevu:
LPM je jednoduchý, ale může predikovat hodnoty mimo interval \([0,1]\), zatímco logit a probit
používají nelineární transformační funkci. Koeficienty logitu se interpretují přes odds ratio, pro
změnu pravděpodobnosti se používají marginální efekty. U neuspořádaných kategorií se používá
multinomický logit nebo probit, u čítacích dat Poissonův model, negativně binomický model a varianty
pro přebytek nul. Prakticky je důležité řešit overdispersion, IIA, expozici, offset a správnou
interpretaci koeficientů.

\subsection{Nejdůležitější termíny}
\begin{terminy}
  \defterm{GLM}{Zobecněný lineární model spojující náhodnou složku, lineární prediktor a link funkci.}
  \defterm{Link funkce}{Funkce propojující střední hodnotu závislé proměnné s lineárním prediktorem.}
  \defterm{LPM}{Lineární pravděpodobnostní model pro binární \(y\); jednoduchý, ale omezený
  predikcemi mimo \([0,1]\) a heteroskedasticitou.}
  \defterm{Logit}{Binární model používající logistickou distribuční funkci pro modelování
  pravděpodobnosti.}
  \defterm{Probit}{Binární model používající distribuční funkci standardního normálního rozdělení.}
  \defterm{Odds}{Poměr pravděpodobnosti jevu a pravděpodobnosti, že jev nenastane.}
  \defterm{Odds ratio}{Násobná změna odds při jednotkové změně vysvětlující proměnné.}
  \defterm{Marginální efekt}{Změna predikované pravděpodobnosti nebo očekávané hodnoty při změně
  vysvětlující proměnné.}
  \defterm{Multinomický logit}{Model pro neuspořádanou kategoriální závislou proměnnou s více než
  dvěma kategoriemi.}
  \defterm{IIA}{Independence of irrelevant alternatives; předpoklad multinomického logitu, že poměr
  pravděpodobností dvou voleb nezávisí na dalších alternativách.}
  \defterm{Poissonova regrese}{Model pro čítací data, kde se logaritmus očekávaného počtu modeluje
  lineárním prediktorem.}
  \defterm{Overdispersion}{Situace, kdy je rozptyl čítací proměnné větší než její střední hodnota.}
  \defterm{Negativně binomický model}{Čítací model vhodný při overdispersion.}
  \defterm{Zero-inflated model}{Model pro čítací data s nadměrným počtem nul.}
  \defterm{Tobit}{Model pro cenzorovanou závislou proměnnou.}
\end{terminy}

\section{Časové řady a panelová data}

\subsection{Shrnutí}
Časové řady řeší závislost pozorování v čase. Před modelováním je nutné řadu vykreslit, posoudit
trend, sezonnost, stacionaritu a autokorelaci a potom volit mezi AR, MA, ARMA, ARIMA, SARIMA,
ADL/ARDL, VAR nebo VECM podle cíle a povahy dat. U predikce je důležitá out-of-sample chyba a kontrola
reziduí, nejen dobrý fit. Panelová data kombinují průřezovou a časovou dimenzi a umožňují kontrolovat
časově neměnnou nepozorovanou heterogenitu. Pooled OLS je jednoduchý, ale často zkreslený; FD a FE
odstraňují individuální efekty, RE vyžaduje silnější předpoklad nekorelovanosti a two-way FE přidává
společné časové šoky.

\subsection{Nejdůležitější termíny}
\begin{terminy}
  \defterm{Časová řada}{Posloupnost pozorování uspořádaných v čase.}
  \defterm{Trend}{Dlouhodobý systematický vývoj úrovně časové řady.}
  \defterm{Sezonnost}{Pravidelně se opakující vzorec v čase, například podle měsíců nebo čtvrtletí.}
  \defterm{Stacionarita}{Vlastnost řady, jejíž základní charakteristiky se v čase nemění.}
  \defterm{Autokorelace}{Korelace hodnot časové řady s jejími zpožděnými hodnotami.}
  \defterm{ACF}{Autokorelační funkce ukazující korelace řady s různými zpožděními.}
  \defterm{PACF}{Parciální autokorelační funkce měřící vztah se zpožděním po očištění o kratší
  zpoždění.}
  \defterm{AR model}{Autoregresní model, kde aktuální hodnota závisí na minulých hodnotách řady.}
  \defterm{MA model}{Model klouzavého průměru, kde aktuální hodnota závisí na minulých šocích.}
  \defterm{ARIMA}{Model kombinující autoregresi, diferencování a klouzavý průměr.}
  \defterm{Jednotkový kořen}{Vlastnost nestacionární řady, u které šoky mají trvalý dopad.}
  \defterm{Kointegrace}{Dlouhodobý rovnovážný vztah mezi nestacionárními řadami.}
  \defterm{VAR}{Vektorový autoregresní model pro více vzájemně provázaných časových řad.}
  \defterm{Panelová data}{Data kombinující pozorování více jednotek v čase.}
  \defterm{Fixed effects}{Panelový odhad odstraňující časově neměnné nepozorované individuální
  charakteristiky.}
  \defterm{Random effects}{Panelový model, který individuální efekt chápe jako náhodnou složku
  nekorelovanou s regresory.}
  \defterm{Hausmanův test}{Test používaný pro volbu mezi fixed effects a random effects.}
\end{terminy}

\section{Strojové učení}

\subsection{Shrnutí}
Strojové učení je přístup, ve kterém se model učí ze vzorových dat a zlepšuje predikce, klasifikace
nebo rozhodování bez ručního naprogramování všech pravidel. Základní dělení zahrnuje učení s učitelem,
učení bez učitele, semi-supervised learning, reinforcement learning, doporučování a detekci anomálií.
Supervised learning řeší hlavně klasifikaci a regresi, unsupervised learning shlukování, asociační
pravidla a redukci dimenze. Typický projekt odpovídá logice CRISP-DM: business porozumění, data,
příprava, modelování, vyhodnocení a nasazení. Model se hodnotí na validačních a testovacích datech
vhodnými metrikami. Klíčová rizika jsou overfitting, data leakage, nevyvážené třídy, špatná metrika
a přehlédnutí nákladů chyb.

\subsection{Nejdůležitější termíny}
\begin{terminy}
  \defterm{Strojové učení}{Oblast AI, kde se model učí vztahy a vzory z dat.}
  \defterm{Feature}{Vstupní atribut nebo prediktor používaný modelem.}
  \defterm{Target}{Cílová proměnná, kterou se model snaží predikovat nebo vysvětlit.}
  \defterm{Učení s učitelem}{Trénování modelu na datech se známou cílovou proměnnou.}
  \defterm{Učení bez učitele}{Hledání struktury v datech bez známé cílové proměnné.}
  \defterm{Klasifikace}{Supervised úloha s kategoriální cílovou proměnnou.}
  \defterm{Regrese}{Supervised úloha s číselnou cílovou proměnnou.}
  \defterm{Reinforcement learning}{Učení agenta, který volí akce v prostředí a učí se z odměn a
  penalizací.}
  \defterm{Trénovací množina}{Data použitá pro odhad parametrů modelu.}
  \defterm{Validační množina}{Data použitá pro ladění modelu a hyperparametrů.}
  \defterm{Testovací množina}{Data použitá pro finální odhad výkonu modelu na nových datech.}
  \defterm{Overfitting}{Přeučení; model se příliš přizpůsobí trénovacím datům a špatně generalizuje.}
  \defterm{Data leakage}{Únik informace z budoucnosti nebo z cílové proměnné do trénovacích příznaků.}
  \defterm{CRISP-DM}{Iterativní metodika datového projektu: business understanding, data understanding,
  data preparation, modeling, evaluation a deployment.}
  \defterm{No-Free-Lunch}{Princip, že neexistuje jeden algoritmus nejlepší pro všechny problémy.}
\end{terminy}

\section{Shlukové a klasifikační úlohy}

\subsection{Shrnutí}
Klasifikace přiřazuje objekty do předem známých tříd a hodnotí se například přes accuracy, precision,
recall, \(F_1\), ROC AUC nebo confusion matrix. Rozhodovací stromy jsou interpretovatelné modely,
které dělí data podle atributů; ID3 používá entropii a information gain, C4.5 přidává gain ratio,
spojité atributy a pruning. Shlukování je učení bez učitele, jehož cílem je najít homogenní a dobře
oddělené skupiny. K-means minimalizuje vzdálenost k centroidům, hierarchické shlukování vytváří
dendrogram a DBSCAN hledá husté oblasti. Asociační pravidla hledají vztahy mezi položkami a hodnotí
se podporou, confidence a liftem; Apriori využívá anti-monotónní vlastnost podpory.

\subsection{Nejdůležitější termíny}
\begin{terminy}
  \defterm{Klasifikace}{Úloha přiřazení objektu do jedné z předem známých tříd.}
  \defterm{Confusion matrix}{Tabulka skutečných a predikovaných tříd používaná pro hodnocení
  klasifikace.}
  \defterm{Precision}{Podíl správně pozitivních predikcí mezi všemi pozitivními predikcemi.}
  \defterm{Recall}{Podíl zachycených pozitivních případů mezi všemi skutečně pozitivními případy.}
  \defterm{\(F_1\) score}{Harmonický průměr precision a recall.}
  \defterm{Rozhodovací strom}{Model tvořený vnitřními rozhodovacími uzly, větvemi a listy s predikcí.}
  \defterm{Entropie}{Míra nečistoty nebo nejistoty tříd v uzlu.}
  \defterm{Information gain}{Snížení entropie po rozdělení uzlu podle atributu.}
  \defterm{Gain ratio}{Normalizovaná varianta informačního zisku omezující preferenci atributů s
  mnoha hodnotami.}
  \defterm{Gini index}{Míra nečistoty uzlu často používaná v algoritmu CART.}
  \defterm{Pruning}{Prořezávání stromu pro omezení přeučení.}
  \defterm{Shlukování}{Unsupervised úloha hledající skupiny podobných objektů.}
  \defterm{Centroid}{Reprezentant shluku, typicky průměr bodů ve shluku.}
  \defterm{K-means}{Shlukovací algoritmus minimalizující součet čtverců vzdáleností bodů od centroidů.}
  \defterm{DBSCAN}{Hustotní shlukovací metoda založená na parametrech \(\varepsilon\) a \(minPts\).}
  \defterm{Asociační pravidlo}{Pravidlo tvaru antecedent \(\Rightarrow\) consequent vyjadřující
  souvýskyt položek.}
  \defterm{Support}{Podpora; podíl transakcí obsahujících danou položkovou množinu nebo pravidlo.}
  \defterm{Confidence}{Spolehlivost pravidla; podmíněná pravděpodobnost consequentu při výskytu
  antecedentu.}
  \defterm{Lift}{Poměr skutečného souvýskytu k souvýskytu očekávanému při nezávislosti.}
\end{terminy}

\section{Ensemblové modelování}

\subsection{Shrnutí}
Ensemblové modely kombinují více dílčích modelů do jednoho výsledného modelu, aby snížily chybu,
zvýšily robustnost nebo přesnost. Bagging vytváří bootstrapové vzorky, trénuje modely nezávisle a
agreguje jejich predikce; typickým příkladem je random forest, který navíc náhodně vybírá atributy
při dělení stromů. Boosting trénuje modely sekvenčně tak, že každý další model opravuje chyby
předchozích. Gradient boosting postupně snižuje ztrátovou funkci ve směru záporného gradientu,
XGBoost je efektivní regularizovaná implementace silná hlavně pro tabulková data. Bagging je obvykle
robustnější a méně citlivý, boosting často přesnější, ale náročnější na ladění a prevenci přeučení.

\subsection{Nejdůležitější termíny}
\begin{terminy}
  \defterm{Ensemble}{Model kombinující predikce více dílčích modelů.}
  \defterm{Hard voting}{Hlasování tříd, kde výsledná třída odpovídá většinové predikci modelů.}
  \defterm{Soft voting}{Agregace predikovaných pravděpodobností tříd.}
  \defterm{Bagging}{Bootstrap aggregating; trénování modelů na bootstrapových vzorcích a agregace
  jejich predikcí.}
  \defterm{Bootstrap sample}{Vzorek vytvořený výběrem s navracením z původních dat.}
  \defterm{Out-of-bag error}{Odhad chyby u baggingu na pozorováních, která nebyla v bootstrapovém
  vzorku daného modelu.}
  \defterm{Random forest}{Ensemble rozhodovacích stromů využívající bagging a náhodný výběr atributů
  při dělení.}
  \defterm{Boosting}{Sekvenční ensemble, kde další slabý model zlepšuje chyby předchozího souboru
  modelů.}
  \defterm{Weak learner}{Jednoduchý model, který je jen o něco lepší než náhodné hádání, ale v
  kombinaci může vytvořit silný model.}
  \defterm{Gradient boosting}{Boosting formulovaný jako postupné minimalizování ztrátové funkce
  přidáváním modelů ve směru záporného gradientu.}
  \defterm{Learning rate}{Koeficient zmenšující příspěvek každého dalšího modelu v boostingu.}
  \defterm{AdaBoost}{Boosting algoritmus zvyšující váhu špatně klasifikovaných pozorování.}
  \defterm{XGBoost}{Efektivní regularizovaná implementace gradient boostingu se stromy.}
  \defterm{Stacking}{Kombinace modelů pomocí meta-modelu, který se učí z jejich predikcí.}
  \defterm{SHAP}{Metoda vysvětlující predikci jako součet příspěvků jednotlivých proměnných.}
\end{terminy}

\section{Mělké neuronové sítě}

\subsection{Shrnutí}
Mělká neuronová síť má vstupní vrstvu, jednu skrytou vrstvu a výstupní vrstvu. Neuron počítá lineární
kombinaci vstupů, přičítá bias a aplikuje aktivační funkci. Díky nelineárním aktivacím síť modeluje
nelineární vztahy. Predikce vzniká dopředným průchodem, chyba se měří ztrátovou funkcí a váhy se
upravují optimalizátorem podle gradientů vypočtených backpropagation. Důležitá je volba architektury,
aktivačních funkcí, výstupní vrstvy, loss function, learning rate, batch velikosti, inicializace vah
a regularizace. Mělké sítě jsou flexibilní, ale méně interpretovatelné, citlivé na škálování a
nastavení hyperparametrů a mohou se přeučit.

\subsection{Nejdůležitější termíny}
\begin{terminy}
  \defterm{Neuron}{Výpočetní jednotka počítající \(a=\phi(w^\top x+b)\).}
  \defterm{Mělká neuronová síť}{Síť s jednou skrytou vrstvou mezi vstupem a výstupem.}
  \defterm{Váha}{Parametr určující sílu vlivu vstupu na neuron.}
  \defterm{Bias}{Aditivní parametr posouvající aktivační hodnotu neuronu.}
  \defterm{Aktivační funkce}{Nelineární funkce umožňující síti modelovat nelineární vztahy.}
  \defterm{Sigmoid}{Aktivační funkce mapující hodnoty do intervalu \((0,1)\), často pro binární
  výstup.}
  \defterm{Tanh}{Aktivační funkce mapující hodnoty do intervalu \((-1,1)\).}
  \defterm{ReLU}{Aktivační funkce \(\max(0,x)\), běžná ve skrytých vrstvách.}
  \defterm{Softmax}{Funkce převádějící vektor skóre na pravděpodobnosti tříd.}
  \defterm{Loss function}{Ztrátová funkce měřící nesoulad mezi predikcí a skutečností.}
  \defterm{Forward propagation}{Dopředný výpočet predikce přes vrstvy sítě.}
  \defterm{Backpropagation}{Efektivní výpočet gradientů ztráty podle vah pomocí řetězového pravidla.}
  \defterm{Gradient descent}{Optimalizační postup aktualizující parametry proti směru gradientu.}
  \defterm{Learning rate}{Velikost kroku při aktualizaci vah.}
  \defterm{Epoch}{Jeden průchod celou trénovací množinou.}
  \defterm{Mini-batch}{Menší dávka pozorování použitá pro jednu aktualizaci vah.}
  \defterm{Early stopping}{Ukončení trénování při zhoršování nebo stagnaci validační chyby.}
\end{terminy}

\section{Hluboké učení}

\subsection{Shrnutí}
Hluboké učení skládá více nelineárních vrstev za sebou a učí se hierarchické reprezentace dat. Je
úspěšné díky velkým datům, GPU, lepším optimalizátorům, regularizaci a specializovaným architekturám.
MLP se používá pro obecné dopředné modelování, CNN pro obraz a lokální strukturu, RNN/LSTM/GRU pro
sekvence, transformery pro jazyk, multimodální data a dlouhé závislosti a autoencodery pro reprezentace,
rekonstrukci a detekci anomálií. Prakticky je důležité nejen trénování, ale také datové pipeline,
validace, MLOps, deployment, monitoring, škálování, bezpečnost a interpretovatelnost. Slabiny hlubokého
učení jsou náročnost na data a výpočet, riziko přeučení, black-box povaha, bias, drift a právní či
etické dopady.

\subsection{Nejdůležitější termíny}
\begin{terminy}
  \defterm{Hluboká neuronová síť}{Neuronová síť s více skrytými vrstvami.}
  \defterm{MLP}{Vícevrstvý perceptron; obecná plně propojená dopředná síť.}
  \defterm{CNN}{Konvoluční neuronová síť vhodná pro data s lokální strukturou, zejména obrazy.}
  \defterm{Konvoluce}{Operace posuvného filtru nad vstupem, která detekuje lokální vzory.}
  \defterm{Pooling}{Zmenšení prostorového rozměru reprezentace, například max poolingem.}
  \defterm{Data augmentation}{Umělé rozšíření trénovacích dat transformacemi, například otočením nebo
  ořezem obrazu.}
  \defterm{RNN}{Rekurentní neuronová síť zpracovávající sekvenční data pomocí skrytého stavu.}
  \defterm{LSTM}{Rekurentní architektura s branami navržená pro lepší zachycení dlouhodobých závislostí.}
  \defterm{GRU}{Jednodušší bránová rekurentní architektura podobná LSTM.}
  \defterm{Embedding}{Hustá vektorová reprezentace slov, položek nebo jiných objektů.}
  \defterm{Transformer}{Architektura založená na attention mechanismu, dobře paralelizovatelná a
  účinná pro jazykové úlohy.}
  \defterm{Self-attention}{Mechanismus, ve kterém si prvky sekvence váženě vybírají relevantní
  informace z ostatních prvků téže sekvence.}
  \defterm{Transfer learning}{Využití předtrénovaného modelu nebo reprezentace pro novou úlohu.}
  \defterm{Autoencoder}{Síť učící se zakódovat vstup do latentní reprezentace a rekonstruovat jej.}
  \defterm{Dropout}{Regularizace náhodným vypínáním neuronů během trénování.}
  \defterm{Batch normalization}{Normalizace aktivací ve vrstvě pro stabilnější a rychlejší trénování.}
  \defterm{MLOps}{Soubor postupů pro správu životního cyklu modelu od vývoje po monitoring v produkci.}
\end{terminy}

\section{Nestrukturovaná a částečně strukturovaná data}

\subsection{Shrnutí}
Nestrukturovaná a semistrukturovaná data je nutné převést do strojově zpracovatelné podoby.
Semistrukturovaná data jako XML, JSON a HTML mají vnitřní dokumentovou nebo stromovou strukturu, ale
ne pevné relační schéma. Nestrukturovaný text se tokenizuje, normalizuje a reprezentuje pomocí termů,
vah nebo embeddingů. Porovnávání vzorů zahrnuje regulární výrazy, edit distance a sémantickou
podobnost. Vyhledávání dokumentů stojí na indexování, zejména inverted indexu, a může používat
booleovský nebo vektorový model. Extrakce informací převádí text na strukturovaná data pomocí NER,
entity linking, relation extraction nebo template filling. Web scraping získává data z webu, ale je
nutné řešit technická omezení, kvalitu, právo, etiku a ochranu osobních údajů.

\subsection{Nejdůležitější termíny}
\begin{terminy}
  \defterm{Strukturovaná data}{Data s pevným schématem, typicky relační tabulky.}
  \defterm{Semistrukturovaná data}{Data s vnitřní strukturou bez pevného tabulkového schématu,
  například XML, JSON nebo HTML.}
  \defterm{Nestrukturovaná data}{Data bez pevné tabulkové struktury, například text, dokumenty,
  obrázky, audio nebo video.}
  \defterm{XML}{Značkovací jazyk se stromovou strukturou, elementy, atributy a možností validace.}
  \defterm{JSON}{Lehký dokumentový formát tvořený objekty, poli a dvojicemi klíč-hodnota.}
  \defterm{CSV}{Jednoduchý textový formát pro tabulková data s oddělovačem sloupců.}
  \defterm{DOM}{Stromová reprezentace HTML dokumentu používaná pro výběr prvků stránky.}
  \defterm{Regulární výraz}{Formální textový vzor pro vyhledávání, validaci, extrakci nebo nahrazování
  částí textu.}
  \defterm{Edit distance}{Míra rozdílu mezi řetězci, často počet operací potřebných k převodu jednoho
  řetězce na druhý.}
  \defterm{Inverted index}{Index mapující termy na dokumenty nebo pozice, ve kterých se vyskytují.}
  \defterm{Booleovský model vyhledávání}{Model pracující s logickými operátory a přesnou shodou
  termů.}
  \defterm{Vektorový model}{Reprezentace dokumentů a dotazů jako vektorů, které se porovnávají
  podobností.}
  \defterm{Kosínová podobnost}{Míra podobnosti vektorů založená na kosinu úhlu mezi nimi.}
  \defterm{NER}{Named entity recognition; rozpoznávání pojmenovaných entit v textu.}
  \defterm{Entity linking}{Propojení zmínky entity v textu s konkrétní entitou ve znalostní bázi.}
  \defterm{XPath}{Jazyk pro výběr uzlů ve stromu XML nebo HTML dokumentu.}
\end{terminy}

\section{Zpracování textových dat}

\subsection{Shrnutí}
Zpracování textových dat převádí přirozený jazyk do podoby vhodné pro vyhledávání, klasifikaci,
modelování, sumarizaci nebo extrakci informací. Základní NLP pipeline zahrnuje tokenizaci, rozdělení
vět, normalizaci, lemmatizaci, POS tagging, syntaktické zpracování, NER, relation extraction a
coreference resolution. Text lze reprezentovat klasicky pomocí bag-of-words, n-gramů a TF-IDF, nebo
moderně pomocí embeddingů a transformerů. Word2Vec učí slovní vektory z kontextu, transformery
vytvářejí kontextové reprezentace pomocí self-attention a RAG kombinuje vyhledávání relevantních
dokumentů s generováním odpovědi. Mezi hlavní úlohy patří klasifikace textu, sentiment analysis,
topic modeling, clustering, sémantické vyhledávání, question answering a sumarizace.

\subsection{Nejdůležitější termíny}
\begin{terminy}
  \defterm{NLP}{Natural language processing; zpracování přirozeného jazyka počítačem.}
  \defterm{Tokenizace}{Rozdělení textu na tokeny, například slova, interpunkci nebo subword jednotky.}
  \defterm{Normalizace textu}{Úprava textu do jednotné podoby, například převod velikosti písmen,
  odstranění šumu nebo sjednocení tvarů.}
  \defterm{Lemmatizace}{Převedení slova na základní slovníkový tvar.}
  \defterm{Stemming}{Hrubší ořez slov na kořen nebo kmen, často bez plné lingvistické správnosti.}
  \defterm{POS tagging}{Přiřazení slovních druhů jednotlivým tokenům.}
  \defterm{Parsing}{Syntaktická analýza věty a vztahů mezi slovy.}
  \defterm{Bag-of-words}{Reprezentace textu jako množiny nebo počtů slov bez pořadí.}
  \defterm{N-gram}{Sekvence \(n\) po sobě jdoucích tokenů.}
  \defterm{TF-IDF}{Váha termu kombinující četnost v dokumentu a vzácnost v korpusu.}
  \defterm{Embedding}{Vektorová reprezentace slova, věty nebo dokumentu v hustém spojitém prostoru.}
  \defterm{Word2Vec}{Metoda učení slovních embeddingů pomocí CBOW nebo Skip-gram úlohy.}
  \defterm{Transformer}{Architektura založená na self-attention, základ moderních jazykových modelů.}
  \defterm{Topic modeling}{Metody hledající latentní témata v kolekci dokumentů, například LDA.}
  \defterm{Sentiment analysis}{Určování postoje, polarity nebo emoce v textu.}
  \defterm{RAG}{Retrieval-augmented generation; kombinace vyhledání relevantních zdrojů a generování
  odpovědi.}
  \defterm{Coreference resolution}{Určování, které výrazy v textu odkazují na stejnou entitu.}
\end{terminy}

\section{Big data}

\subsection{Shrnutí}
Big data jsou data, která kvůli objemu, rychlosti, rozmanitosti nebo nárokům na hodnotu a kvalitu
vyžadují škálovatelné uložení, zpracování a správu. Základní princip spočívá v rozdělení dat na
partitions, paralelním výpočtu a odolnosti proti výpadkům. Hadoop přinesl HDFS a MapReduce, zatímco
Spark nabízí rychlejší a pohodlnější zpracování přes RDD, DataFrame, Spark SQL, MLlib a Structured
Streaming. Streamová data se řeší pomocí Kafky, Spark Structured Streaming nebo Flinku, s důrazem na
event time, okna, stav, checkpointy a offsety. Pro ukládání se používají data lake, lakehouse,
Parquet, partitioning a katalogy metadat. Vedle techniky je nutné řešit governance, kvalitu, bezpečnost,
monitoring, orchestraci, náklady a fyzickou organizaci dat.

\subsection{Nejdůležitější termíny}
\begin{terminy}
  \defterm{Big data}{Data vyžadující škálovatelné technologie pro uložení, zpracování a správu.}
  \defterm{5V}{Volume, velocity, variety, veracity a value jako charakteristiky velkých dat.}
  \defterm{Partition}{Část dat umožňující paralelní zpracování a efektivnější čtení.}
  \defterm{Hadoop}{Ekosystém pro distribuované uložení a dávkové zpracování dat.}
  \defterm{HDFS}{Distribuovaný souborový systém Hadoopu dělící soubory na bloky ukládané na DataNode.}
  \defterm{NameNode}{Uzel HDFS spravující metadata o blocích a umístění souborů.}
  \defterm{MapReduce}{Dávkový výpočetní model s fázemi map, shuffle/sort a reduce.}
  \defterm{Apache Spark}{Distribuovaný výpočetní engine pro dávkové, interaktivní, ML a streamové
  zpracování.}
  \defterm{RDD}{Resilient distributed dataset; základní distribuovaná kolekce ve Sparku.}
  \defterm{DataFrame}{Tabulková distribuovaná datová struktura ve Sparku s optimalizací dotazů.}
  \defterm{Lazy evaluation}{Odložené vyhodnocování operací až do okamžiku akce.}
  \defterm{Apache Kafka}{Distribuovaná event streaming platforma pro publikování, ukládání a čtení
  událostí.}
  \defterm{Data lake}{Úložiště velkého množství dat v surovější podobě, často v objektovém úložišti.}
  \defterm{Lakehouse}{Architektura kombinující flexibilitu data lake s vlastnostmi datového skladu.}
  \defterm{Parquet}{Sloupcový souborový formát vhodný pro analytické čtení.}
  \defterm{CAP teorém}{Při síťové partition musí distribuovaný systém volit kompromis mezi konzistencí
  a dostupností.}
  \defterm{Medallion architektura}{Vrstvení dat typicky na bronze, silver a gold podle kvality a
  zpracovanosti.}
\end{terminy}

\section{Business intelligence}

\subsection{Shrnutí}
Business intelligence převádí podniková data na informace použitelné pro řízení. Typická architektura
zahrnuje zdrojové systémy, ETL nebo ELT, staging, ODS, datový sklad nebo lakehouse, sémantickou vrstvu
a reportovací nástroje. Datový sklad je subjektově orientované, integrované, časově rozlišené a
relativně neměnné úložiště pro rozhodování. Pro analytické modelování je klíčové rozlišit fakta a
dimenze. Hvězdicové schéma je jednoduché a výkonné, sněhové schéma normalizuje dimenze. OLAP umožňuje
multidimenzionální analýzu přes slice, dice, drill-down, roll-up a pivot. Úspěšné BI nestojí pouze na
nástroji, ale na kvalitních datech, jasných definicích metrik, governance, bezpečnosti, monitoringu a
vazbě na skutečné rozhodování.

\subsection{Nejdůležitější termíny}
\begin{terminy}
  \defterm{Business intelligence}{Metody, procesy, datové struktury a nástroje převádějící podniková
  data na informace pro řízení.}
  \defterm{Zdrojový systém}{Interní nebo externí systém, ze kterého BI čerpá data.}
  \defterm{Datový sklad}{Centrální analytické úložiště integrovaných historických dat pro podporu
  rozhodování.}
  \defterm{Datové tržiště}{Menší analytické úložiště zaměřené na konkrétní proces, oddělení nebo
  doménu.}
  \defterm{Sémantická vrstva}{Vrstva převádějící technický datový model na business pojmy, metriky,
  dimenze a hierarchie.}
  \defterm{Faktová tabulka}{Tabulka měřitelných událostí nebo stavů, například prodeje, objednávky
  nebo stavy zásob.}
  \defterm{Dimenzní tabulka}{Tabulka popisující kontext faktů, například čas, produkt, zákazník nebo
  region.}
  \defterm{Granularita faktů}{Úroveň detailu jednoho řádku faktové tabulky.}
  \defterm{Hvězdicové schéma}{Dimenzionální model s centrální faktovou tabulkou a denormalizovanými
  dimenzemi.}
  \defterm{Sněhové schéma}{Dimenzionální model s normalizovanými dimenzemi.}
  \defterm{OLAP}{Analytické zpracování nad multidimenzionálními daty.}
  \defterm{KPI}{Klíčový ukazatel výkonnosti navázaný na cíl a rozhodování.}
  \defterm{Dashboard}{Prezentační vrstva BI s ukazateli, grafy, filtry a interakcemi.}
  \defterm{Data governance}{Pravidla, role a procesy pro správu kvality, odpovědnosti, bezpečnosti
  a využívání dat.}
\end{terminy}

\section{Řízení datového projektu}

\subsection{Shrnutí}
Řízení datového projektu spojuje projektové řízení, analytiku, data engineering, governance a zavedení
výsledků do praxe. Základní logiku dobře vystihuje CRISP-DM: porozumění businessu, porozumění datům,
příprava dat, modelování, evaluace a nasazení. Proces je iterativní, protože cíle, data a modely se
během práce zpřesňují. Vedle CRISP-DM je dobré znát SEMMA, TDSP, agilní řízení, Scrum, DAMA-DMBOK a
MLOps. Projekt potřebuje jasné role, backlog, dokumentaci, verzování, orchestrace pipeline,
experiment tracking, model registry a monitoring. Hlavní rizika jsou nejasný cíl, nekvalitní data,
právní omezení, špatná evaluace, technický dluh a nepřijetí uživateli.

\subsection{Nejdůležitější termíny}
\begin{terminy}
  \defterm{Datový projekt}{Projekt, jehož výstupem je datový produkt, analytické řešení, model,
  dashboard, pipeline nebo rozhodovací podpora.}
  \defterm{Trojimperativ}{Vazba rozsahu, času a nákladů; změna jedné složky ovlivňuje ostatní.}
  \defterm{CRISP-DM}{Iterativní metodika datového projektu s fázemi business understanding, data
  understanding, data preparation, modeling, evaluation a deployment.}
  \defterm{SEMMA}{Metodika sample, explore, modify, model a assess.}
  \defterm{TDSP}{Team Data Science Process; rámec pro týmovou realizaci datově-vědních projektů.}
  \defterm{Scrum}{Agilní rámec založený na sprintech, backlogu, rolích a pravidelné inspekci
  postupu.}
  \defterm{Stakeholder}{Osoba nebo skupina ovlivněná projektem nebo schopná projekt ovlivnit.}
  \defterm{Business vlastník}{Role odpovědná za hodnotu, cíle a praktické využití výstupu.}
  \defterm{Data owner}{Role odpovědná za data, jejich dostupnost, kvalitu, pravidla a oprávnění.}
  \defterm{Backlog}{Seznam požadavků, úkolů a priorit projektu.}
  \defterm{Experiment tracking}{Evidence běhů modelování, parametrů, metrik a artefaktů.}
  \defterm{Model registry}{Správa verzí modelů, jejich stavů, metadat a nasazení.}
  \defterm{Deployment}{Nasazení analytického řešení nebo modelu do používání.}
  \defterm{Monitoring modelu}{Sledování výkonu, driftu, chyb, dostupnosti a dopadů modelu po nasazení.}
  \defterm{Technický dluh}{Budoucí náklady způsobené zkratkami, provizorii a nedostatečnou údržbou
  řešení.}
\end{terminy}

\section{Optimalizační modely}

\subsection{Shrnutí}
Optimalizační model vybírá nejlepší přípustné rozhodnutí vzhledem k účelové funkci a omezením. Lineární
programování má lineární účelovou funkci a lineární omezení; přípustná množina je konvexní a optimum,
pokud existuje, leží ve vrcholu. Simplexová metoda postupuje mezi bazickými přípustnými řešeními a
zlepšuje účelovou funkci. Dualita umožňuje ekonomicky interpretovat duální proměnné jako stínové ceny.
ILP a MILP přidávají celočíselné a binární proměnné pro výběry, fixní náklady, logické podmínky a
kombinatorické struktury. LP relaxace poskytuje mez, branch-and-bound větví a ořezává neperspektivní
větve a cutting planes přidávají platná omezení. Heuristiky a metaheuristiky hledají dobrá řešení bez
záruky optima, což je důležité u velkých a výpočetně těžkých úloh.

\subsection{Nejdůležitější termíny}
\begin{terminy}
  \defterm{Optimalizační model}{Matematický model s rozhodovacími proměnnými, účelovou funkcí a
  omezeními.}
  \defterm{Rozhodovací proměnná}{Proměnná reprezentující volbu, kterou model optimalizuje.}
  \defterm{Účelová funkce}{Kritérium optimalizace, například maximalizace zisku nebo minimalizace
  nákladů.}
  \defterm{Omezení}{Podmínka vymezující přípustná rozhodnutí.}
  \defterm{Lineární programování}{Optimalizace s lineární účelovou funkcí a lineárními omezeními.}
  \defterm{Přípustná množina}{Množina všech řešení splňujících omezení modelu.}
  \defterm{Bazické přípustné řešení}{Řešení odpovídající bázi lineárně nezávislých sloupců a splňující
  omezení včetně nezápornosti.}
  \defterm{Simplexová metoda}{Algoritmus řešení LP pohybující se mezi vrcholy přípustné množiny.}
  \defterm{Duální úloha}{Optimalizační úloha přiřazená primalní úloze, často s ekonomickou interpretací
  omezených zdrojů.}
  \defterm{Stínová cena}{Duální proměnná vyjadřující změnu optimální hodnoty při malém navýšení
  pravé strany omezení.}
  \defterm{MILP}{Mixed-integer linear programming; lineární model se spojitými i celočíselnými nebo
  binárními proměnnými.}
  \defterm{LP relaxace}{Model vzniklý odstraněním celočíselných omezení z ILP nebo MILP.}
  \defterm{Branch-and-bound}{Metoda řešení celočíselných úloh pomocí větvení, LP relaxací a ořezávání
  větví.}
  \defterm{Cutting plane}{Platné omezení přidané k odstranění neceločíselného relaxovaného řešení bez
  ztráty celočíselných řešení.}
  \defterm{TSP}{Travelling salesman problem; úloha nalezení nejkratší okružní trasy přes daná místa.}
  \defterm{Metaheuristika}{Obecný heuristický rámec pro hledání dobrých řešení složitých úloh, například
  simulované žíhání, genetické algoritmy nebo tabu search.}
  \defterm{Optimality gap}{Rozdíl mezi aktuálním nejlepším řešením a mezí optimální hodnoty.}
\end{terminy}

\end{document}
